% =====================================================================
%  MANUEL DE REVISION CAPET SII IE — SESSION 2026
%  Assemblage entrelace : cours (fiches/) + exercices (exercices/),
%  corriges regroupes en annexe.
%
%  COMPILATION (depuis manuel/, 3 passes pour sommaire + renvois) :
%    pdflatex -interaction=nonstopmode main.tex   (x3)
%
%  Ce fichier n'est qu'un ASSEMBLAGE : il ne duplique aucun contenu.
%  Chaque fiche et chaque banque d'exercices reste sa propre source,
%  compilable seule exactement comme avant.
%
%  \subimport (package import) permet aux sources incluses de garder
%  leurs \input relatifs (\input{preamble} cote fiches,
%  % preamble.tex — packages et commandes communes a toutes les fiches CAPET SII IE
\input{glyphtounicode}
\pdfgentounicode=1
\usepackage{cmap}
\usepackage[T1]{fontenc}
\usepackage[utf8]{inputenc}
\usepackage[french]{babel}
% babel[french] active a \begin{document} des raccourcis de ponctuation (:;!?)
% qui cassent la syntaxe to[...] de circuitikz (erreur "+ or - expected") ;
% on les desactive, ce qui ne change que l'espacement fin avant ces signes.
\AtBeginDocument{\shorthandoff{:;!?}}
\usepackage[a4paper, top=1.6cm, bottom=1.6cm, left=1.4cm, right=1.4cm]{geometry}
\setlength{\headheight}{14.5pt}
\usepackage{amsmath, amssymb, amsfonts}
\usepackage{siunitx}
\sisetup{locale=FR, output-decimal-marker={,}, inter-unit-product=\ensuremath{{}\cdot{}},
         group-minimum-digits=4}
% Unites non fournies par siunitx, declarees ici une fois pour toutes (fiches + exercices).
% Les gardes \ifdefined evitent tout conflit avec une redeclaration locale restee dans un
% fichier : ne pas retirer les gardes, ne pas redeclarer ces unites dans les fichiers.
\DeclareSIUnit{\jour}{j}
\ifdefined\var\else\DeclareSIUnit{\var}{var}\fi   % puissance reactive
\ifdefined\vad\else\DeclareSIUnit{\vad}{VAD}\fi   % puissance deformante
\ifdefined\tr\else\DeclareSIUnit{\tr}{tr}\fi      % tour (tr/min)
\usepackage{array, booktabs, multirow, tabularx}
\usepackage{xcolor}
\usepackage[most, breakable]{tcolorbox}
\usepackage{multicol}
\usepackage{enumitem}
\usepackage{mathtools}
\usepackage{lmodern}
% Refonte visuelle 2026 (direction "moderne coloré") : Lato pour le texte,
% maths sans-serif assorties via newtxsf (charge apres amsmath/amssymb).
\usepackage[default]{lato}
\usepackage{newtxsf}
\usepackage{microtype}
\usepackage{fancyhdr}
\usepackage{tikz}
\usetikzlibrary{positioning,arrows.meta}
\usepackage{circuitikz}
\usepackage{lastpage}

% ── Couleurs par domaine (palette modernisée 2026) ────────────────────
\definecolor{colElec}{RGB}{29, 78, 216}        % Électrotechnique — bleu vif
\definecolor{colElecLight}{RGB}{219, 234, 254} % fond bleu clair
\definecolor{colAuto}{RGB}{21, 128, 61}        % Automatique — vert
\definecolor{colMeca}{RGB}{194, 65, 12}        % Mécanique — orange
\definecolor{colTherm}{RGB}{185, 28, 28}       % Thermique — rouge
\definecolor{colQT}{RGB}{232, 244, 233}        % (hérité, plus utilisé par questiontype)
\definecolor{colJury}{RGB}{255, 248, 225}      % fond jury (crème)
\definecolor{colJuryFrame}{RGB}{202, 138, 4}   % bordure/bandeau jury (ambre)
\definecolor{colGrayLight}{RGB}{245, 245, 247} % fond gris neutre
\definecolor{colSeparator}{RGB}{180, 180, 180} % gris séparateur
\definecolor{colExo}{RGB}{8, 145, 178}          % Banque d'exercices — cyan (distinct des fiches)
% ── Encre et gris de la refonte ───────────────────────────────────────
\definecolor{colInk}{RGB}{23, 32, 42}          % texte courant (quasi-noir)
\definecolor{colSlate}{RGB}{83, 98, 118}       % gris ardoise (méta, en-têtes)
\definecolor{colRule}{RGB}{203, 213, 225}      % filets clairs
\definecolor{colJuryBg}{RGB}{255, 248, 225}    % fond crème jury
\definecolor{colPresence}{RGB}{83, 98, 118}    % bloc 5 (présence dans les sujets)
\AtBeginDocument{\color{colInk}}

% ── Style encadré exercice + corrigé (banque d'exercices) ─────────────
\tcbset{
  exoBox/.style={
    breakable,
    enhanced,
    colback=colExo!8!white,
    colframe=colExo,
    colbacktitle=colExo,
    coltitle=white,
    fonttitle=\bfseries\small,
    left=4pt, right=4pt, top=3pt, bottom=3pt,
    before skip=5pt, after skip=5pt,
  }
}

% ── Bascule enonce / corrige (banque d'exercices) ─────────────────────
% \CAPETCORRIGE est defini en ligne de commande (pdflatex "\def\CAPETCORRIGE{1}\input{...}")
% pour compiler la variante corriges d'un meme fichier source ; par defaut (non defini),
% la compilation produit la variante enonces. \CAPETSHOW{enonce}{corrige} choisit la branche.
%
% TROIS MODES (refonte 2026-09-20, manuel de revision) :
%   - aucun drapeau   -> imprime l'enonce            (variante enonces, INCHANGE)
%   - \CAPETCORRIGE   -> imprime le corrige          (variante corriges, INCHANGE)
%   - \CAPETMANUEL    -> imprime l'enonce ET met le corrige de cote pour l'annexe
%
% Le mode manuel stocke chaque corrige dans une macro (\csgdef), sans ecriture
% fichier ni re-tokenisation : le contenu (circuitikz, accents, %) traverse intact.
% L'annexe les rejoue ensuite dans l'ordre via \CAPETANNEXE.
\usepackage{etoolbox}

\newcounter{CAPETexo}          % numero continu d'exercice dans le manuel
\setcounter{CAPETexo}{0}
\newcounter{CAPETannexe}       % compteur de relecture en annexe

% Titre courant de la notion en cours (pose par \CAPETNOTION), sert d'intertitre annexe
\newcommand{\CAPETnotioncode}{}
\newcommand{\CAPETnotiontitre}{}

% \CAPETNOTION{code}{titre} : ouvre une notion dans le manuel (memorise le libelle,
% et marque le rang du premier exercice de la notion pour l'intertitre d'annexe).
\newcommand{\CAPETNOTION}[2]{%
  \renewcommand{\CAPETnotioncode}{#1}%
  \renewcommand{\CAPETnotiontitre}{#2}%
  \csgdef{CAPETnotionstart\the\numexpr\value{CAPETexo}+1\relax}{#1 -- #2}%
}

% ── \degree robuste ───────────────────────────────────────────────────
% Certains fichiers d'exercices utilisent \degree. Ni le preambule ni siunitx
% ne le definissent, et \providecommand local ne suffit pas : sous hyperref,
% \degree existe deja comme commande factice (\ERROR), donc \providecommand
% s'abstient et l'appel echoue a la compilation du manuel (erreur invisible
% en compilation isolee). On le (re)definit donc inconditionnellement ici,
% apres le chargement des packages, pour tout le corpus.
\AtBeginDocument{\providecommand{\degree}{}\renewcommand{\degree}{\ensuremath{^\circ}}}

% Libelle decoratif : les fichiers d'exercices utilisent aussi \CAPETSHOW dans
% leur \fancyhead et leur bloc de titre, pour afficher "(enonces)" ou "(corriges)".
% Ces appels-la ne sont PAS des exercices : en mode manuel ils ne doivent ni
% consommer un numero, ni alimenter l'annexe. On les reconnait a leur premier
% argument, strictement egal au libelle ci-dessous.
\newcommand{\CAPETLIBELLEENONCES}{(\'enonc\'es)}

\ifdefined\CAPETMANUEL
  % --- Mode manuel : enonce en place, corrige differe vers l'annexe ---
  \newcommand{\CAPETSHOW}[2]{%
    \def\CAPETargun{#1}%
    \ifx\CAPETargun\CAPETLIBELLEENONCES
      #1%            appel decoratif : rendu tel quel, hors numerotation
    \else
      \CAPETSHOWEXO{#1}{#2}%
    \fi
  }
  \newcommand{\CAPETSHOWEXO}[2]{%
    \stepcounter{CAPETexo}%
    \csgdef{CAPETcorrige\theCAPETexo}{#2}%
    \edef\CAPETtmpnum{\theCAPETexo}%
    #1%
    \par\vspace{2pt}%
    {\footnotesize\itshape\color{colExo}%
     \hyperref[capet:cor:\CAPETtmpnum]{Corrige de l'exercice \CAPETtmpnum\ $\rightarrow$ p.~\pageref{capet:cor:\CAPETtmpnum}}%
     \label{capet:eno:\CAPETtmpnum}}%
  }
  % \CAPETANNEXE : rejoue tous les corriges captures, dans l'ordre, avec intertitres
  \newcommand{\CAPETANNEXE}{%
    \setcounter{CAPETannexe}{0}%
    \loop\ifnum\value{CAPETannexe}<\value{CAPETexo}%
      \stepcounter{CAPETannexe}%
      \edef\CAPETcurnum{\theCAPETannexe}%
      \ifcsname CAPETnotionstart\CAPETcurnum\endcsname
        \par\medskip
        {\normalsize\bfseries\color{colExo}\csuse{CAPETnotionstart\CAPETcurnum}}%
        \par\nobreak\smallskip
      \fi
      \begin{tcolorbox}[exoBox, breakable,
        title={Corrige de l'exercice \CAPETcurnum\
               \hfill\normalfont\footnotesize
               \hyperref[capet:eno:\CAPETcurnum]{enonce p.~\pageref{capet:eno:\CAPETcurnum}}}]%
      \label{capet:cor:\CAPETcurnum}%
      \csuse{CAPETcorrige\CAPETcurnum}%
      \end{tcolorbox}%
    \repeat
  }
\else
  % --- Modes historiques, comportement strictement inchange ---
  \newcommand{\CAPETSHOW}[2]{\ifdefined\CAPETCORRIGE #2\else #1\fi}
  \newcommand{\CAPETANNEXE}{}
\fi

% ── Difficulté 1 à 3 (étoiles pleines, amssymb) ───────────────────────
\newcommand{\difficulte}[1]{%
  \ifcase#1\relax
  \or $\bigstar$\;niveau 1/3%
  \or $\bigstar\bigstar$\;niveau 2/3%
  \or $\bigstar\bigstar\bigstar$\;niveau 3/3%
  \fi
}

% ── Style encadré extrait de document technique (banque d'exercices) ──
\tcbset{
  docTech/.style={
    breakable,
    enhanced,
    colback=colGrayLight,
    colframe=colSeparator,
    fonttitle=\bfseries\small\itshape,
    coltitle=black,
    colbacktitle=colGrayLight,
    left=4pt, right=4pt, top=3pt, bottom=3pt,
    before skip=4pt, after skip=4pt,
  }
}

% ── Style tcolorbox de section principale ────────────────────────────
% Refonte 2026 : bandeau titre coloré arrondi + fond très pâle, sans cadre.
\tcbset{
  ficheSection/.style={
    breakable, enhanced jigsaw,
    frame hidden, boxrule=0pt, arc=1.8mm,
    colback=#1!4!white,
    colbacktitle=#1,
    coltitle=white,
    fonttitle=\sffamily\bfseries,
    left=7pt, right=7pt, top=5pt, bottom=5pt,
    toptitle=3pt, bottomtitle=3pt,
    before skip=10pt, after skip=8pt,
  }
}

% ── Style encadré question type ──────────────────────────────────────
% Carte blanche, bandeau sombre, liseré très clair.
\tcbset{
  questiontype/.style={
    breakable, enhanced jigsaw,
    colback=white,
    frame hidden, boxrule=0pt, arc=1.8mm,
    colbacktitle=colElec!78!black,
    coltitle=white,
    fonttitle=\sffamily\bfseries\small,
    left=7pt, right=7pt, top=4pt, bottom=4pt,
    toptitle=2.5pt, bottomtitle=2.5pt,
    before skip=8pt, after skip=8pt,
    borderline={0.6pt}{0pt}{colElec!12!white},
  }
}

% ── Style encadré jury ───────────────────────────────────────────────
% Bandeau ambre, fond crème, sans cadre.
\tcbset{
  jury/.style={
    breakable, enhanced jigsaw,
    colback=colJuryBg,
    frame hidden, boxrule=0pt, arc=1.8mm,
    colbacktitle=colJuryFrame,
    coltitle=white,
    fonttitle=\sffamily\bfseries,
    left=7pt, right=7pt, top=4pt, bottom=4pt,
    toptitle=3pt, bottomtitle=3pt,
    before skip=10pt, after skip=8pt,
  }
}

% ── En-tête / pied de page communs (refonte 2026) ─────────────────────
% #1 code fiche, #2 titre court, #3 couleur domaine
\newcommand{\ficheHeaders}[3]{%
  \fancyhead[L]{\small\sffamily{\bfseries\textcolor{#3}{#1}}\hspace{0.7em}\textcolor{colSlate}{#2}}%
  \fancyhead[R]{\small\sffamily\textcolor{colSlate}{CAPET SII IE\, --- \,Session 2026}}%
  \fancyfoot[C]{\small\sffamily\textcolor{colSlate}{\thepage\,/\,\pageref{LastPage}}}%
  \fancyhead[C]{}\fancyfoot[L]{}\fancyfoot[R]{}%
}

% ── Bloc en-tête de fiche (BLOC 0, refonte 2026) ──────────────────────
% #1 code  #2 titre  #3 couleur domaine  #4 domaine  #5 ligne présence  #6 sous-thèmes
\newcommand{\ficheEntete}[6]{%
  \begin{tcolorbox}[enhanced, breakable,
    frame hidden, boxrule=0pt, arc=2.5mm,
    interior style={left color=#3, right color=#3!75!black},
    left=9pt, right=9pt, top=7pt, bottom=7pt,
    before skip=0pt, after skip=6pt]
    {\color{white}\sffamily\bfseries\LARGE #1\quad #2}\\[4pt]
    {\color{#3!22!white}\sffamily\small \MakeUppercase{#4}\quad\textbar\quad #5}\\[1pt]
    {\color{#3!35!white}\sffamily\footnotesize Sous-thèmes : #6}
  \end{tcolorbox}%
}

% ── Mise en page ─────────────────────────────────────────────────────
\pagestyle{fancy}
\fancyhf{}
\renewcommand{\headrulewidth}{0.5pt}
\renewcommand{\footrulewidth}{0pt}
\AtBeginDocument{%
  \renewcommand{\headrule}{{\color{colRule}\hrule width\headwidth height 0.5pt}}%
}

\setlist[itemize]{leftmargin=1.2em, itemsep=1pt, topsep=2pt, parsep=0pt}
\setlist[enumerate]{leftmargin=1.2em, itemsep=3pt, topsep=2pt, parsep=0pt}

\setlength{\parskip}{2pt}
\setlength{\columnsep}{10pt}
 cote exercices).
% =====================================================================
\documentclass[11pt,a4paper,twoside]{article}

\def\CAPETMANUEL{1}                 % mode manuel : corriges differes vers l'annexe
\newcommand{\CAPETINCLUDEMODE}{1}   % les sources n'ouvrent pas leur propre document

\usepackage{import}
% preamble.tex — packages et commandes communes a toutes les fiches CAPET SII IE
\input{glyphtounicode}
\pdfgentounicode=1
\usepackage{cmap}
\usepackage[T1]{fontenc}
\usepackage[utf8]{inputenc}
\usepackage[french]{babel}
% babel[french] active a \begin{document} des raccourcis de ponctuation (:;!?)
% qui cassent la syntaxe to[...] de circuitikz (erreur "+ or - expected") ;
% on les desactive, ce qui ne change que l'espacement fin avant ces signes.
\AtBeginDocument{\shorthandoff{:;!?}}
\usepackage[a4paper, top=1.6cm, bottom=1.6cm, left=1.4cm, right=1.4cm]{geometry}
\setlength{\headheight}{14.5pt}
\usepackage{amsmath, amssymb, amsfonts}
\usepackage{siunitx}
\sisetup{locale=FR, output-decimal-marker={,}, inter-unit-product=\ensuremath{{}\cdot{}},
         group-minimum-digits=4}
% Unites non fournies par siunitx, declarees ici une fois pour toutes (fiches + exercices).
% Les gardes \ifdefined evitent tout conflit avec une redeclaration locale restee dans un
% fichier : ne pas retirer les gardes, ne pas redeclarer ces unites dans les fichiers.
\DeclareSIUnit{\jour}{j}
\ifdefined\var\else\DeclareSIUnit{\var}{var}\fi   % puissance reactive
\ifdefined\vad\else\DeclareSIUnit{\vad}{VAD}\fi   % puissance deformante
\ifdefined\tr\else\DeclareSIUnit{\tr}{tr}\fi      % tour (tr/min)
\usepackage{array, booktabs, multirow, tabularx}
\usepackage{xcolor}
\usepackage[most, breakable]{tcolorbox}
\usepackage{multicol}
\usepackage{enumitem}
\usepackage{mathtools}
\usepackage{lmodern}
% Refonte visuelle 2026 (direction "moderne coloré") : Lato pour le texte,
% maths sans-serif assorties via newtxsf (charge apres amsmath/amssymb).
\usepackage[default]{lato}
\usepackage{newtxsf}
\usepackage{microtype}
\usepackage{fancyhdr}
\usepackage{tikz}
\usetikzlibrary{positioning,arrows.meta}
\usepackage{circuitikz}
\usepackage{lastpage}

% ── Couleurs par domaine (palette modernisée 2026) ────────────────────
\definecolor{colElec}{RGB}{29, 78, 216}        % Électrotechnique — bleu vif
\definecolor{colElecLight}{RGB}{219, 234, 254} % fond bleu clair
\definecolor{colAuto}{RGB}{21, 128, 61}        % Automatique — vert
\definecolor{colMeca}{RGB}{194, 65, 12}        % Mécanique — orange
\definecolor{colTherm}{RGB}{185, 28, 28}       % Thermique — rouge
\definecolor{colQT}{RGB}{232, 244, 233}        % (hérité, plus utilisé par questiontype)
\definecolor{colJury}{RGB}{255, 248, 225}      % fond jury (crème)
\definecolor{colJuryFrame}{RGB}{202, 138, 4}   % bordure/bandeau jury (ambre)
\definecolor{colGrayLight}{RGB}{245, 245, 247} % fond gris neutre
\definecolor{colSeparator}{RGB}{180, 180, 180} % gris séparateur
\definecolor{colExo}{RGB}{8, 145, 178}          % Banque d'exercices — cyan (distinct des fiches)
% ── Encre et gris de la refonte ───────────────────────────────────────
\definecolor{colInk}{RGB}{23, 32, 42}          % texte courant (quasi-noir)
\definecolor{colSlate}{RGB}{83, 98, 118}       % gris ardoise (méta, en-têtes)
\definecolor{colRule}{RGB}{203, 213, 225}      % filets clairs
\definecolor{colJuryBg}{RGB}{255, 248, 225}    % fond crème jury
\definecolor{colPresence}{RGB}{83, 98, 118}    % bloc 5 (présence dans les sujets)
\AtBeginDocument{\color{colInk}}

% ── Style encadré exercice + corrigé (banque d'exercices) ─────────────
\tcbset{
  exoBox/.style={
    breakable,
    enhanced,
    colback=colExo!8!white,
    colframe=colExo,
    colbacktitle=colExo,
    coltitle=white,
    fonttitle=\bfseries\small,
    left=4pt, right=4pt, top=3pt, bottom=3pt,
    before skip=5pt, after skip=5pt,
  }
}

% ── Bascule enonce / corrige (banque d'exercices) ─────────────────────
% \CAPETCORRIGE est defini en ligne de commande (pdflatex "\def\CAPETCORRIGE{1}\input{...}")
% pour compiler la variante corriges d'un meme fichier source ; par defaut (non defini),
% la compilation produit la variante enonces. \CAPETSHOW{enonce}{corrige} choisit la branche.
%
% TROIS MODES (refonte 2026-09-20, manuel de revision) :
%   - aucun drapeau   -> imprime l'enonce            (variante enonces, INCHANGE)
%   - \CAPETCORRIGE   -> imprime le corrige          (variante corriges, INCHANGE)
%   - \CAPETMANUEL    -> imprime l'enonce ET met le corrige de cote pour l'annexe
%
% Le mode manuel stocke chaque corrige dans une macro (\csgdef), sans ecriture
% fichier ni re-tokenisation : le contenu (circuitikz, accents, %) traverse intact.
% L'annexe les rejoue ensuite dans l'ordre via \CAPETANNEXE.
\usepackage{etoolbox}

\newcounter{CAPETexo}          % numero continu d'exercice dans le manuel
\setcounter{CAPETexo}{0}
\newcounter{CAPETannexe}       % compteur de relecture en annexe

% Titre courant de la notion en cours (pose par \CAPETNOTION), sert d'intertitre annexe
\newcommand{\CAPETnotioncode}{}
\newcommand{\CAPETnotiontitre}{}

% \CAPETNOTION{code}{titre} : ouvre une notion dans le manuel (memorise le libelle,
% et marque le rang du premier exercice de la notion pour l'intertitre d'annexe).
\newcommand{\CAPETNOTION}[2]{%
  \renewcommand{\CAPETnotioncode}{#1}%
  \renewcommand{\CAPETnotiontitre}{#2}%
  \csgdef{CAPETnotionstart\the\numexpr\value{CAPETexo}+1\relax}{#1 -- #2}%
}

% ── \degree robuste ───────────────────────────────────────────────────
% Certains fichiers d'exercices utilisent \degree. Ni le preambule ni siunitx
% ne le definissent, et \providecommand local ne suffit pas : sous hyperref,
% \degree existe deja comme commande factice (\ERROR), donc \providecommand
% s'abstient et l'appel echoue a la compilation du manuel (erreur invisible
% en compilation isolee). On le (re)definit donc inconditionnellement ici,
% apres le chargement des packages, pour tout le corpus.
\AtBeginDocument{\providecommand{\degree}{}\renewcommand{\degree}{\ensuremath{^\circ}}}

% Libelle decoratif : les fichiers d'exercices utilisent aussi \CAPETSHOW dans
% leur \fancyhead et leur bloc de titre, pour afficher "(enonces)" ou "(corriges)".
% Ces appels-la ne sont PAS des exercices : en mode manuel ils ne doivent ni
% consommer un numero, ni alimenter l'annexe. On les reconnait a leur premier
% argument, strictement egal au libelle ci-dessous.
\newcommand{\CAPETLIBELLEENONCES}{(\'enonc\'es)}

\ifdefined\CAPETMANUEL
  % --- Mode manuel : enonce en place, corrige differe vers l'annexe ---
  \newcommand{\CAPETSHOW}[2]{%
    \def\CAPETargun{#1}%
    \ifx\CAPETargun\CAPETLIBELLEENONCES
      #1%            appel decoratif : rendu tel quel, hors numerotation
    \else
      \CAPETSHOWEXO{#1}{#2}%
    \fi
  }
  \newcommand{\CAPETSHOWEXO}[2]{%
    \stepcounter{CAPETexo}%
    \csgdef{CAPETcorrige\theCAPETexo}{#2}%
    \edef\CAPETtmpnum{\theCAPETexo}%
    #1%
    \par\vspace{2pt}%
    {\footnotesize\itshape\color{colExo}%
     \hyperref[capet:cor:\CAPETtmpnum]{Corrige de l'exercice \CAPETtmpnum\ $\rightarrow$ p.~\pageref{capet:cor:\CAPETtmpnum}}%
     \label{capet:eno:\CAPETtmpnum}}%
  }
  % \CAPETANNEXE : rejoue tous les corriges captures, dans l'ordre, avec intertitres
  \newcommand{\CAPETANNEXE}{%
    \setcounter{CAPETannexe}{0}%
    \loop\ifnum\value{CAPETannexe}<\value{CAPETexo}%
      \stepcounter{CAPETannexe}%
      \edef\CAPETcurnum{\theCAPETannexe}%
      \ifcsname CAPETnotionstart\CAPETcurnum\endcsname
        \par\medskip
        {\normalsize\bfseries\color{colExo}\csuse{CAPETnotionstart\CAPETcurnum}}%
        \par\nobreak\smallskip
      \fi
      \begin{tcolorbox}[exoBox, breakable,
        title={Corrige de l'exercice \CAPETcurnum\
               \hfill\normalfont\footnotesize
               \hyperref[capet:eno:\CAPETcurnum]{enonce p.~\pageref{capet:eno:\CAPETcurnum}}}]%
      \label{capet:cor:\CAPETcurnum}%
      \csuse{CAPETcorrige\CAPETcurnum}%
      \end{tcolorbox}%
    \repeat
  }
\else
  % --- Modes historiques, comportement strictement inchange ---
  \newcommand{\CAPETSHOW}[2]{\ifdefined\CAPETCORRIGE #2\else #1\fi}
  \newcommand{\CAPETANNEXE}{}
\fi

% ── Difficulté 1 à 3 (étoiles pleines, amssymb) ───────────────────────
\newcommand{\difficulte}[1]{%
  \ifcase#1\relax
  \or $\bigstar$\;niveau 1/3%
  \or $\bigstar\bigstar$\;niveau 2/3%
  \or $\bigstar\bigstar\bigstar$\;niveau 3/3%
  \fi
}

% ── Style encadré extrait de document technique (banque d'exercices) ──
\tcbset{
  docTech/.style={
    breakable,
    enhanced,
    colback=colGrayLight,
    colframe=colSeparator,
    fonttitle=\bfseries\small\itshape,
    coltitle=black,
    colbacktitle=colGrayLight,
    left=4pt, right=4pt, top=3pt, bottom=3pt,
    before skip=4pt, after skip=4pt,
  }
}

% ── Style tcolorbox de section principale ────────────────────────────
% Refonte 2026 : bandeau titre coloré arrondi + fond très pâle, sans cadre.
\tcbset{
  ficheSection/.style={
    breakable, enhanced jigsaw,
    frame hidden, boxrule=0pt, arc=1.8mm,
    colback=#1!4!white,
    colbacktitle=#1,
    coltitle=white,
    fonttitle=\sffamily\bfseries,
    left=7pt, right=7pt, top=5pt, bottom=5pt,
    toptitle=3pt, bottomtitle=3pt,
    before skip=10pt, after skip=8pt,
  }
}

% ── Style encadré question type ──────────────────────────────────────
% Carte blanche, bandeau sombre, liseré très clair.
\tcbset{
  questiontype/.style={
    breakable, enhanced jigsaw,
    colback=white,
    frame hidden, boxrule=0pt, arc=1.8mm,
    colbacktitle=colElec!78!black,
    coltitle=white,
    fonttitle=\sffamily\bfseries\small,
    left=7pt, right=7pt, top=4pt, bottom=4pt,
    toptitle=2.5pt, bottomtitle=2.5pt,
    before skip=8pt, after skip=8pt,
    borderline={0.6pt}{0pt}{colElec!12!white},
  }
}

% ── Style encadré jury ───────────────────────────────────────────────
% Bandeau ambre, fond crème, sans cadre.
\tcbset{
  jury/.style={
    breakable, enhanced jigsaw,
    colback=colJuryBg,
    frame hidden, boxrule=0pt, arc=1.8mm,
    colbacktitle=colJuryFrame,
    coltitle=white,
    fonttitle=\sffamily\bfseries,
    left=7pt, right=7pt, top=4pt, bottom=4pt,
    toptitle=3pt, bottomtitle=3pt,
    before skip=10pt, after skip=8pt,
  }
}

% ── En-tête / pied de page communs (refonte 2026) ─────────────────────
% #1 code fiche, #2 titre court, #3 couleur domaine
\newcommand{\ficheHeaders}[3]{%
  \fancyhead[L]{\small\sffamily{\bfseries\textcolor{#3}{#1}}\hspace{0.7em}\textcolor{colSlate}{#2}}%
  \fancyhead[R]{\small\sffamily\textcolor{colSlate}{CAPET SII IE\, --- \,Session 2026}}%
  \fancyfoot[C]{\small\sffamily\textcolor{colSlate}{\thepage\,/\,\pageref{LastPage}}}%
  \fancyhead[C]{}\fancyfoot[L]{}\fancyfoot[R]{}%
}

% ── Bloc en-tête de fiche (BLOC 0, refonte 2026) ──────────────────────
% #1 code  #2 titre  #3 couleur domaine  #4 domaine  #5 ligne présence  #6 sous-thèmes
\newcommand{\ficheEntete}[6]{%
  \begin{tcolorbox}[enhanced, breakable,
    frame hidden, boxrule=0pt, arc=2.5mm,
    interior style={left color=#3, right color=#3!75!black},
    left=9pt, right=9pt, top=7pt, bottom=7pt,
    before skip=0pt, after skip=6pt]
    {\color{white}\sffamily\bfseries\LARGE #1\quad #2}\\[4pt]
    {\color{#3!22!white}\sffamily\small \MakeUppercase{#4}\quad\textbar\quad #5}\\[1pt]
    {\color{#3!35!white}\sffamily\footnotesize Sous-thèmes : #6}
  \end{tcolorbox}%
}

% ── Mise en page ─────────────────────────────────────────────────────
\pagestyle{fancy}
\fancyhf{}
\renewcommand{\headrulewidth}{0.5pt}
\renewcommand{\footrulewidth}{0pt}
\AtBeginDocument{%
  \renewcommand{\headrule}{{\color{colRule}\hrule width\headwidth height 0.5pt}}%
}

\setlist[itemize]{leftmargin=1.2em, itemsep=1pt, topsep=2pt, parsep=0pt}
\setlist[enumerate]{leftmargin=1.2em, itemsep=3pt, topsep=2pt, parsep=0pt}

\setlength{\parskip}{2pt}
\setlength{\columnsep}{10pt}


% --- Mise en page RECTO VERSO (reliure) -------------------------------
% Surcharge locale : le preambule partage est en marges symetriques 1,4 cm,
% adapte a une fiche imprimee seule. Pour un ouvrage relie il faut une marge
% INTERIEURE plus large (cote pliure). \geometry surcharge sans toucher
% fiches/preamble.tex, donc les fiches autonomes gardent leur mise en page.
\geometry{inner=2.2cm, outer=1.4cm, top=1.6cm, bottom=1.6cm}

% En-tetes alternes : le titre de notion vers l'exterieur de chaque page.
% Redefinition de \ficheHeaders (defini dans le preambule partage) pour que
% les fiches et exercices inclus adoptent le comportement recto verso.
\renewcommand{\ficheHeaders}[3]{%
  \fancyhf{}%
  % Titre de notion vers l'EXTERIEUR de la double page, mention generale a l'interieur.
  % Verso (page paire) : exterieur = gauche. Recto (page impaire) : exterieur = droite.
  \fancyhead[LE]{\small\sffamily{\bfseries\textcolor{#3}{#1}}\hspace{0.7em}\textcolor{colSlate}{#2}}%
  \fancyhead[RE]{\small\sffamily\textcolor{colSlate}{CAPET SII IE\, --- \,Session 2026}}%
  \fancyhead[RO]{\small\sffamily{\bfseries\textcolor{#3}{#1}}\hspace{0.7em}\textcolor{colSlate}{#2}}%
  \fancyhead[LO]{\small\sffamily\textcolor{colSlate}{CAPET SII IE\, --- \,Session 2026}}%
  \fancyfoot[C]{\small\sffamily\textcolor{colSlate}{\thepage\,/\,\pageref{LastPage}}}%
}

\usepackage[hidelinks, bookmarks=true, bookmarksnumbered=true]{hyperref}
\hypersetup{
  pdftitle={Manuel de revision CAPET SII IE - Session 2026},
  pdfsubject={Preparation au CAPET SII option Ingenierie electrique},
  pdfcreator={LaTeX},
}
\usepackage{bookmark}

% parties_rectoverso.tex — intercalaires (variante RECTO VERSO) de partie du manuel de revision.
% \CAPETPARTIE{numero}{titre}{sous-titre}{couleur} : page de garde de partie.
% Les sauts de page vivent ICI et dans main.tex, JAMAIS dans les fiches
% (regle CLAUDE.md : aucun \newpage/\clearpage dans fiches/T*.tex).

\newcommand{\CAPETPARTIE}[4]{%
  \cleardoublepage
  \thispagestyle{empty}%
  \null\vspace*{3.2cm}%
  \begin{center}
    \begin{tcolorbox}[enhanced, width=\textwidth,
      colback=#4, colframe=#4, arc=2.5mm,
      left=14pt, right=14pt, top=16pt, bottom=16pt]
      \centering
      {\color{white}\sffamily\bfseries\large PARTIE #1}\\[8pt]
      {\color{white}\sffamily\bfseries\LARGE #2}
    \end{tcolorbox}
    \vspace{1.6cm}

    \begin{minipage}{0.82\textwidth}
      \centering
      {\color{colSlate}\sffamily\large #3}
    \end{minipage}
  \end{center}
  \vfill
  \cleardoublepage
}


% --- Ouverture d'une notion dans le manuel -----------------------------
% #1 code, #2 titre, #3 couleur : saut de page + entree de sommaire +
% memorisation de la notion pour les intertitres d'annexe.
\newcommand{\CAPETOUVRENOTION}[3]{%
  \cleardoublepage
  \phantomsection
  \addcontentsline{toc}{subsection}{\texorpdfstring{\textcolor{#3}{#1}}{#1} --- #2}%
  \CAPETNOTION{#1}{#2}%
}

% --- Bandeau d'ouverture de la section exercices d'une notion ----------
\newcommand{\CAPETBANDEAUEXOS}[2]{%
  \par\medskip
  \begin{tcolorbox}[enhanced, breakable, colback=colExo, colframe=colExo,
    left=6pt, right=6pt, top=4pt, bottom=4pt, arc=1.5mm,
    before skip=8pt, after skip=6pt]
    {\color{white}\bfseries\large Exercices --- #1\quad #2}
  \end{tcolorbox}
}

\begin{document}
\pagenumbering{gobble}

% =====================================================================
%  COUVERTURE
% =====================================================================
\begin{titlepage}
\thispagestyle{empty}
\begin{tikzpicture}[remember picture,overlay]
  \fill[colGrayLight] (current page.south west) rectangle (current page.north east);
  \fill[colElec] (current page.south west) rectangle ([xshift=3.8cm]current page.north west);
  \fill[colAuto] ([xshift=3.8cm,yshift=-4.2cm]current page.north west) rectangle ([xshift=5.1cm]current page.south west);
  \fill[colTherm] ([xshift=5.1cm,yshift=-8.2cm]current page.north west) rectangle ([xshift=6.3cm]current page.south west);
  \fill[colMeca] ([xshift=6.3cm,yshift=-11.2cm]current page.north west) rectangle ([xshift=7.2cm]current page.south west);
  \draw[colElec!35, line width=1.2pt] ([xshift=8.2cm,yshift=-2.0cm]current page.north west) rectangle ([xshift=-1.2cm,yshift=2.0cm]current page.south east);
\end{tikzpicture}

\vspace*{1.8cm}
\hspace*{4.9cm}
\begin{minipage}{0.68\textwidth}
  {\color{colElec}\Large\bfseries CAPET SII}\\[3pt]
  {\color{black!70}\large Option Ing\'enierie \'electrique}\\[1.2cm]
  {\fontsize{28}{32}\selectfont\bfseries Manuel de r\'evision}\\[0.25cm]
  {\fontsize{17}{21}\selectfont\bfseries Cours, exercices et corrig\'es}\\[0.8cm]
  {\large Session 2026}\\[0.35cm]
  {\normalsize 30 notions : \'electrotechnique, \'electronique, automatique,
   m\'ecanique appliqu\'ee et thermique. Chaque notion associe la fiche de
   synth\`ese et ses exercices d'entra\^inement ; les corrig\'es sont
   regroup\'es en annexe.}
\end{minipage}

\vfill

\hspace*{4.9cm}
\begin{minipage}{0.68\textwidth}
\begin{tcolorbox}[enhanced, colback=white, colframe=colElec, boxrule=0.8pt,
  arc=2mm, left=8pt, right=8pt, top=8pt, bottom=8pt]
\textbf{Mode d'emploi}\\[4pt]
Chaque notion se lit dans l'ordre : \textbf{la fiche de cours}, puis
\textbf{les exercices}. Travaillez chaque exercice avant d'ouvrir son
corrig\'e : le renvoi en fin d'\'enonc\'e donne sa page en annexe.
\end{tcolorbox}
{\large\bfseries CAPET SII IE}\\[2pt]
{\normalsize Base documentaire \LaTeX{} pour impression}
\end{minipage}
\end{titlepage}

% =====================================================================
%  SOMMAIRE
% =====================================================================
\clearpage
\pdfbookmark[1]{Sommaire}{toc}
{\renewcommand{\contentsname}{Sommaire}
 \tableofcontents}

% Le liminaire (couverture + sommaire) doit occuper un nombre PAIR de pages pour
% que la page 1 arabe tombe sur un recto (page physique impaire). Sinon la parite
% logique et la parite physique sont decalees et \cleardoublepage ouvre les
% notions a gauche. \cleardoublepage ici complete le liminaire si besoin.
\cleardoublepage
\pagenumbering{arabic}
\setcounter{page}{1}

% =====================================================================
%  PARTIE I — ELECTROTECHNIQUE (T01 a T17)
% =====================================================================
\phantomsection\addcontentsline{toc}{section}{Partie I --- \'Electrotechnique}
\CAPETPARTIE{I}{\'Electrotechnique}{Grandeurs, puissance, conversion statique, machines, \'energie et distribution --- notions T01 \`a T17}{colElec}

\CAPETOUVRENOTION{T01}{Grandeurs \'electriques et \'electrocin\'etique}{colElec}
\subimport{../fiches/}{T01_Grandeurs_Electrocinetique}
\CAPETBANDEAUEXOS{T01}{Grandeurs \'electriques et \'electrocin\'etique}
\subimport{../exercices/}{T01_exercices}

\CAPETOUVRENOTION{T02}{Puissance en r\'egime alternatif sinuso\"idal}{colElec}
\subimport{../fiches/}{T02_Puissance_AC}
\CAPETBANDEAUEXOS{T02}{Puissance en r\'egime alternatif sinuso\"idal}
\subimport{../exercices/}{T02_exercices}

\CAPETOUVRENOTION{T03}{Repr\'esentation complexe, imp\'edances}{colElec}
\subimport{../fiches/}{T03_Representations_Complexes_Impedances}
\CAPETBANDEAUEXOS{T03}{Repr\'esentation complexe, imp\'edances}
\subimport{../exercices/}{T03_exercices}

\CAPETOUVRENOTION{T04}{Harmoniques r\'eseau, THD}{colElec}
\subimport{../fiches/}{T04_Harmoniques_THD}
\CAPETBANDEAUEXOS{T04}{Harmoniques r\'eseau, THD}
\subimport{../exercices/}{T04_exercices}

\CAPETOUVRENOTION{T05}{Ph\'enom\`enes thermiques}{colTherm}
\subimport{../fiches/}{T05_Phenomenes_Thermiques}
\CAPETBANDEAUEXOS{T05}{Ph\'enom\`enes thermiques}
\subimport{../exercices/}{T05_exercices}

\CAPETOUVRENOTION{T06}{Transformateur}{colElec}
\subimport{../fiches/}{T06_Transformateur}
\CAPETBANDEAUEXOS{T06}{Transformateur}
\subimport{../exercices/}{T06_exercices}

\CAPETOUVRENOTION{T07}{Hacheurs}{colElec}
\subimport{../fiches/}{T07_Hacheurs}
\CAPETBANDEAUEXOS{T07}{Hacheurs}
\subimport{../exercices/}{T07_exercices}

\CAPETOUVRENOTION{T08}{Redresseurs}{colElec}
\subimport{../fiches/}{T08_Redresseurs}
\CAPETBANDEAUEXOS{T08}{Redresseurs}
\subimport{../exercices/}{T08_exercices}

\CAPETOUVRENOTION{T09}{Onduleurs}{colElec}
\subimport{../fiches/}{T09_Onduleurs}
\CAPETBANDEAUEXOS{T09}{Onduleurs}
\subimport{../exercices/}{T09_exercices}

\CAPETOUVRENOTION{T10}{M\'ecanique et transmissions}{colMeca}
\subimport{../fiches/}{T10_Mecanique}
\CAPETBANDEAUEXOS{T10}{M\'ecanique et transmissions}
\subimport{../exercices/}{T10_exercices}

\CAPETOUVRENOTION{T11}{Machine \`a courant continu}{colElec}
\subimport{../fiches/}{T11_Machine_CC}
\CAPETBANDEAUEXOS{T11}{Machine \`a courant continu}
\subimport{../exercices/}{T11_exercices}

\CAPETOUVRENOTION{T12}{Machine synchrone et brushless}{colElec}
\subimport{../fiches/}{T12_Machine_Synchrone}
\CAPETBANDEAUEXOS{T12}{Machine synchrone et brushless}
\subimport{../exercices/}{T12_exercices}

\CAPETOUVRENOTION{T13}{Machine asynchrone}{colElec}
\subimport{../fiches/}{T13_Machine_Asynchrone}
\CAPETBANDEAUEXOS{T13}{Machine asynchrone}
\subimport{../exercices/}{T13_exercices}

\CAPETOUVRENOTION{T14}{Variation de vitesse et de couple}{colElec}
\subimport{../fiches/}{T14_Variation_Vitesse}
\CAPETBANDEAUEXOS{T14}{Variation de vitesse et de couple}
\subimport{../exercices/}{T14_exercices}

\CAPETOUVRENOTION{T15}{Stockage d'\'energie \'electrique}{colElec}
\subimport{../fiches/}{T15_Stockage_Energie}
\CAPETBANDEAUEXOS{T15}{Stockage d'\'energie \'electrique}
\subimport{../exercices/}{T15_exercices}

\CAPETOUVRENOTION{T16}{\'Energies renouvelables}{colElec}
\subimport{../fiches/}{T16_Energies_Renouvelables}
\CAPETBANDEAUEXOS{T16}{\'Energies renouvelables}
\subimport{../exercices/}{T16_exercices}

\CAPETOUVRENOTION{T17}{Distribution et s\'ecurit\'e \'electrique}{colElec}
\subimport{../fiches/}{T17_Distribution_Securite}
\CAPETBANDEAUEXOS{T17}{Distribution et s\'ecurit\'e \'electrique}
\subimport{../exercices/}{T17_exercices}

% =====================================================================
%  PARTIE II — ELECTRONIQUE (T18 a T25)
% =====================================================================
\phantomsection\addcontentsline{toc}{section}{Partie II --- \'Electronique}
\CAPETPARTIE{II}{\'Electronique}{Composants, amplification, filtrage, num\'erisation, capteurs, embarqu\'e et transmission --- notions T18 \`a T25}{colElec}

\CAPETOUVRENOTION{T18}{Composants \'electroniques}{colElec}
\subimport{../fiches/}{T18_Composants_Electroniques}
\CAPETBANDEAUEXOS{T18}{Composants \'electroniques}
\subimport{../exercices/}{T18_exercices}

\CAPETOUVRENOTION{T19}{Amplificateur op\'erationnel}{colElec}
\subimport{../fiches/}{T19_AOP}
\CAPETBANDEAUEXOS{T19}{Amplificateur op\'erationnel}
\subimport{../exercices/}{T19_exercices}

\CAPETOUVRENOTION{T20}{R\'egime sinuso\"idal et filtres}{colElec}
\subimport{../fiches/}{T20_Filtres_RC_RL}
\CAPETBANDEAUEXOS{T20}{R\'egime sinuso\"idal et filtres}
\subimport{../exercices/}{T20_exercices}

\CAPETOUVRENOTION{T21}{Fonctions de transfert, Bode}{colElec}
\subimport{../fiches/}{T21_Fonctions_Transfert_Bode}
\CAPETBANDEAUEXOS{T21}{Fonctions de transfert, Bode}
\subimport{../exercices/}{T21_exercices}

\CAPETOUVRENOTION{T22}{CAN/CNA et num\'erisation}{colElec}
\subimport{../fiches/}{T22_CAN_CNA_Numerisation}
\CAPETBANDEAUEXOS{T22}{CAN/CNA et num\'erisation}
\subimport{../exercices/}{T22_exercices}

\CAPETOUVRENOTION{T23}{Capteurs et conditionnement}{colElec}
\subimport{../fiches/}{T23_Capteurs_Conditionnement}
\CAPETBANDEAUEXOS{T23}{Capteurs et conditionnement}
\subimport{../exercices/}{T23_exercices}

\CAPETOUVRENOTION{T24}{Syst\`emes embarqu\'es}{colElec}
\subimport{../fiches/}{T24_Systemes_Embarques}
\CAPETBANDEAUEXOS{T24}{Syst\`emes embarqu\'es}
\subimport{../exercices/}{T24_exercices}

\CAPETOUVRENOTION{T25}{Transmission num\'erique}{colElec}
\subimport{../fiches/}{T25_Transmission_Numerique}
\CAPETBANDEAUEXOS{T25}{Transmission num\'erique}
\subimport{../exercices/}{T25_exercices}

% =====================================================================
%  PARTIE III — AUTOMATIQUE (T26 a T29)
% =====================================================================
\phantomsection\addcontentsline{toc}{section}{Partie III --- Automatique}
\CAPETPARTIE{III}{Automatique}{Mod\'elisation, asservissement, correcteurs et logique --- notions T26 \`a T29}{colAuto}

\CAPETOUVRENOTION{T26}{Mod\'elisation des syst\`emes lin\'eaires}{colAuto}
\subimport{../fiches/}{T26_Modelisation_Systemes}
\CAPETBANDEAUEXOS{T26}{Mod\'elisation des syst\`emes lin\'eaires}
\subimport{../exercices/}{T26_exercices}

\CAPETOUVRENOTION{T27}{Syst\`emes boucl\'es et correcteurs}{colAuto}
\subimport{../fiches/}{T27_Systemes_Boucles_Correcteurs}
\CAPETBANDEAUEXOS{T27}{Syst\`emes boucl\'es et correcteurs}
\subimport{../exercices/}{T27_exercices}

\CAPETOUVRENOTION{T28}{Logique combinatoire}{colAuto}
\subimport{../fiches/}{T28_Logique_Combinatoire}
\CAPETBANDEAUEXOS{T28}{Logique combinatoire}
\subimport{../exercices/}{T28_exercices}

\CAPETOUVRENOTION{T29}{Logique s\'equentielle, GRAFCET}{colAuto}
\subimport{../fiches/}{T29_Logique_Sequentielle_GRAFCET}
\CAPETBANDEAUEXOS{T29}{Logique s\'equentielle, GRAFCET}
\subimport{../exercices/}{T29_exercices}

% =====================================================================
%  PARTIE IV — MESURAGE ET INSTRUMENTATION (T30)
% =====================================================================
\phantomsection\addcontentsline{toc}{section}{Partie IV --- Mesurage et instrumentation}
\CAPETPARTIE{IV}{Mesurage et instrumentation}{Cha\^ine de mesure, incertitudes et sch\'ematisation normalis\'ee --- notion T30}{colElec}

\CAPETOUVRENOTION{T30}{Mesurage et instrumentation}{colElec}
\subimport{../fiches/}{T30_Mesurage_Instrumentation}
\CAPETBANDEAUEXOS{T30}{Mesurage et instrumentation}
\subimport{../exercices/}{T30_exercices}

% =====================================================================
%  ANNEXE A — CORRIGES
% =====================================================================
\phantomsection\addcontentsline{toc}{section}{Annexe A --- Corrig\'es des exercices}
\CAPETPARTIE{A}{Corrig\'es des exercices}{Corrig\'es d\'etaill\'es de tous les exercices du manuel, dans l'ordre des notions}{colExo}

\ficheHeaders{Annexe A}{Corrig\'es des exercices}{colExo}
\CAPETANNEXE

\end{document}
